\section{Motivation}
As outlined above, the MLton default array-of-pointers is disasterous from the
perspective of efficiently utilizing the memory controller, in terms of both
bandwidth and latency. The MLton priorities in the MLton approach appear to be
inverted: rather than treating array-of-pointers elimination as the
highest-priority among the potential flattening options, it is treated as a
final, best-effort transformation. For common array-oriented workloads (HPC,
graphics, etc), this has demonstrably-poor results. These workloads appear to be
under-represented in the MLton benchmark suite, but they have good coverage in
parallel-ml-bench, as we'll see below.

How can we generalize the argument-centric \texttt{Flatten} transformation to
arrays? In the \texttt{SSA} representation, the action of the \texttt{Flatten}
transformation is to dissovle the association between the conjoined elements by
providing them separate names, i.e.;
\begin{minted}{python}
  t: a * b * c * ... -->
  x: a, y: b, z: c, ...
\end{minted}
Mapping this on to arrays, we would, likewise, provide a name for each element
in the tuple, i.e.:
\begin{minted}{python}
  arr: (a * b * c * ...) array -->
  arr_x: a array, arr_y: b array, arr_z: c array * ...
\end{minted}
Considering the \texttt{Flatten} transformation on the \texttt{Machine}
representation (focusing on the non-tail-recursive case), the action is to copy
the \hl{conjoined elements} themselves into the buffer carrying on the function
arguments (rather than a pointer to them), i.e.:
\begin{minted}{cpp}
  *stack.Push<void*>() = t ->
  *stack.Push<typeA>() = t[0], 
  *stack.Push<typeB>() = t[1],
   ...
\end{minted}
The analogue for arrays is to copy the elements themselves into the array, i.e.:
\begin{minted}{cpp}
  arr[i] = t ->
  arr[i * tuplelen(*t)] = t[0],
  arr[i * tuplelen(*t) + 1] = t[1],
  ...
\end{minted}
Speaking somewhat loosely, these two choices of strategies correspond to the
familiar struct-of-arrays (SoA) vs array-of-structs (AoS) layouts. The AoS
layout is the natural default for a C++ programmer, i.e.:
\begin{minted}{cpp}
  struct Point {
    int x;
    int y;
  };
  std::vector<Point> array_of_structs;
\end{minted}
SoA is the natural default for SIMD operations, since it makes it simple to load
contiguous chunks of data into SIMD registers:
\begin{minted}{cpp}
  struct PointSoA {
    std::vector<int> xs;
    std::vector<int> ys;
  };
\end{minted}
Libraries and frameworks to support SIMD CPU programming commonly offer support
for accessing SoA-format data in an AoS-like format, e.g. Intel's ISPC,
introduced in \citet{ispc-paper}. ``Manual'' SoA layouts are also a standard
technique in high-performance GPU programming, as discussed in e.g.
\citet{cuda-textbook}.

The case of MLton differs slightly from the AoS/SoA tradeoff in that, as
described above, our default memory layout is an array-of-pointers:
\begin{minted}{cpp}
  struct Point {
    int x;
    int y;
  };
  std::vector<Point*> array_of_pointers;
\end{minted}
which lacks a name in standard performance-engineering parlance, since it is
virtually never the right choice to achieve high \hl{performance}. We will refer
to it here as AoP when necessary.

In order to handle the case where the consuming operations of the array elements
cannot or should not be flattened, \texttt{ShallowFlatten} inserts ``pack''
operations for elements read from the AoS/SoA versions of the containers, and
``unpack'' operations for elements written to the AoS/SoA versions. This lets us
preserve the pre-existing names and types in the consuming code (for the
elements), simplifying the transformation at the cost of some
potentially-unnecessary heap operations if the intermediate tuples can't be
optimzied away by later passes. While this can in theory result in a net
performance loss, choosing the right heuristic for flattening is a question for
the benchmarks, and the high cost of the AoP layout means that the bar to clear
is very low.

For the containers themselves, however, we need to update the types everywhere
that they appear: not only in variable bindings and function arguments, but also
in \texttt{datatype} declarations, including \texttt{datatype}s introduced as
part of closure conversion. This makes it difficult to support partial versions
of the transformation in order to enable cases where some, but not all AoP
containers in the program are \texttt{ShallowFlatten}ed: the implementation
below assumes all-or-nothing conversion.


%% deep flatten would be great but we observe that it often doesn't work in
%% practice. the priority seems to be inverted: deepflatten thinks of container
%% flattening as a minor best-effort win, but really it's crucial for HPC
%% workloads. these workloads seem to be under-represented in the MLton
%% benchmark suite, but they appear in parallel-ml-bench

%% what are the natural generalizatons of flattening? one strategy is that we
%% pull the pointers out of the containers and drop the association between
%% elements. this seems like the natural generalization of tuple flattening.
%% another strategy is that we push the elements inside the container and drop
%% the tuple structure. this is closer to what deep flatten does

%% these two strategies coincide with AoS vs SoA. SoA is the preferred format
%% for SIMD. it is a little odd to call the other version AoS (C++ programmers
%% would think of AoS as "the default" rather than something you can opt-into),
%% but it helps to make the distinction clear

%% in order to preserve the nearby types, we insert pack/unpack operations. this
%% allows us to be more aggressive than deep flatten. by running before the
%% other flattening passes (in SSA IR rather than SSA2) we create opportunities
%% for those tuples to be eliminated anyway (though this choice introduces some
%% fixable limitations -- specifically about uniform types). it's possible in
%% theory to imagine that the more-aggressive choice is worse if it requires us
%% to do a lot of extra pack/unpack, but that's a question for the benchmarks

\hl{TODO: FIGURE FROM GEMINI -- NEED TO INTEGRATE INTO THE TEXT}

Figure~\ref{fig:memory-layouts} illustrates the memory layout differences between
the default Array-of-Pointers (AoP) representation and the two shallow flattening
strategies, Array-of-Structs (AoS) and Structure-of-Arrays (SoA), alongside their
representative C++ container definitions.

\begin{figure}[htbp]
  \centering
  \tikzset{
    cell/.style={draw=black!75, line width=0.6pt, minimum height=0.62cm, font=\small, align=center},
    idx/.style={font=\scriptsize\color{textMuted}},
    labelnode/.style={font=\small\bfseries, anchor=east},
    ptr/.style={circle, fill=colorPtrBorder, inner sep=1.8pt},
    arrptr/.style={-{Stealth[length=4.5pt, width=3.5pt]}, line width=0.8pt, color=colorPtrBorder},
    brace/.style={decorate, decoration={brace, amplitude=3.5pt, mirror}, line width=0.6pt, black!70}
  }

  % -----------------------------------------------------------
  % (a) Array-of-Pointers (AoP)
  % -----------------------------------------------------------
  \begin{subfigure}{\textwidth}
    \centering
    \begin{minipage}[c]{0.43\textwidth}
\begin{minted}[fontsize=\footnotesize,frame=single,framesep=2.5mm,rulecolor=\color{black!20}]{cpp}
struct Point {
  int x;
  int y;
};
std::vector<Point*> array_of_pointers;
\end{minted}
    \end{minipage}%
    \hfill
    \begin{minipage}[c]{0.55\textwidth}
      \centering
      \begin{tikzpicture}
        % Pointer array
        \node[labelnode] at (-0.2, 0) {\texttt{arr}:};
        \foreach \i in {0,1,2,3} {
          \node[cell, fill=colorPtr, draw=colorPtrBorder, minimum width=1.3cm] (ptr\i) at (\i*1.3 + 0.65, 0) {};
          \node[idx] at ($(ptr\i.north) + (0, 0.18)$) {[\i]};
          \node[ptr] (dot\i) at (ptr\i.center) {};
        }

        % Heap Tuples
        \node[labelnode] at (-0.2, -1.45) {Heap:};
        \foreach \i/\addr in {0/0x1040, 1/0x28A0, 2/0x14C0, 3/0x3F00} {
          \pgfmathsetmacro{\cx}{\i*1.3 + 0.65}
          \node[cell, fill=colorHdr, draw=colorHdrBorder, minimum width=0.34cm, minimum height=0.52cm, font=\tiny\color{textMuted}] 
            (h\i) at (\cx - 0.38, -1.45) {hdr};
          \node[cell, fill=colorX, draw=colorXBorder, minimum width=0.38cm, minimum height=0.52cm, right=0cm of h\i] 
            (hx\i) {\scriptsize$x_\i$};
          \node[cell, fill=colorY, draw=colorYBorder, minimum width=0.38cm, minimum height=0.52cm, right=0cm of hx\i] 
            (hy\i) {\scriptsize$y_\i$};
          \node[draw=black!25, dashed, rounded corners=2pt, fit=(h\i)(hy\i), inner sep=1pt] (box\i) {};
          \node[font=\tiny\ttfamily\color{textMuted}, below=1pt of box\i] {@\addr};

          % Pointer arrow
          \draw[arrptr] (dot\i) -- (box\i.north);
        }
      \end{tikzpicture}
    \end{minipage}
    \caption{Array-of-Pointers (AoP): Pointer array to scattered heap-allocated tuples.}
    \label{fig:layout-aop}
  \end{subfigure}

  \vspace{1.2em}

  % -----------------------------------------------------------
  % (b) Array-of-Structs (AoS)
  % -----------------------------------------------------------
  \begin{subfigure}{\textwidth}
    \centering
    \begin{minipage}[c]{0.43\textwidth}
\begin{minted}[fontsize=\footnotesize,frame=single,framesep=2.5mm,rulecolor=\color{black!20}]{cpp}
struct Point {
  int x;
  int y;
};
std::vector<Point> array_of_structs;
\end{minted}
    \end{minipage}%
    \hfill
    \begin{minipage}[c]{0.55\textwidth}
      \centering
      \begin{tikzpicture}
        \node[labelnode] at (-0.2, 0) {\texttt{arr}:};
        \foreach \i in {0,1,2,3} {
          \pgfmathsetmacro{\cx}{\i*1.3 + 0.325}
          \node[cell, fill=colorX, draw=colorXBorder, minimum width=0.65cm] (ax\i) at (\cx, 0) {\scriptsize$x_\i$};
          \node[cell, fill=colorY, draw=colorYBorder, minimum width=0.65cm, right=0cm of ax\i] (ay\i) {\scriptsize$y_\i$};

          \draw[brace] ($(ax\i.south west) + (0.04,-0.08)$) -- ($(ay\i.south east) + (-0.04,-0.08)$)
            node[midway, below=2pt, font=\scriptsize] {\texttt{arr[\i]}};

          \node[idx] at ($(ax\i.north) + (0, 0.18)$) {[\pgfmathparse{int(\i*2)}\pgfmathresult]};
          \node[idx] at ($(ay\i.north) + (0, 0.18)$) {[\pgfmathparse{int(\i*2+1)}\pgfmathresult]};
        }
      \end{tikzpicture}
    \end{minipage}
    \caption{Array-of-Structs (AoS): Single buffer with contiguous interleaved fields.}
    \label{fig:layout-aos}
  \end{subfigure}

  \vspace{1.2em}

  % -----------------------------------------------------------
  % (c) Structure-of-Arrays (SoA)
  % -----------------------------------------------------------
  \begin{subfigure}{\textwidth}
    \centering
    \begin{minipage}[c]{0.43\textwidth}
\begin{minted}[fontsize=\footnotesize,frame=single,framesep=2.5mm,rulecolor=\color{black!20}]{cpp}
struct PointSoA {
  std::vector<int> xs;
  std::vector<int> ys;
};
PointSoA struct_of_arrays;
\end{minted}
    \end{minipage}%
    \hfill
    \begin{minipage}[c]{0.55\textwidth}
      \centering
      \begin{tikzpicture}
        % xs
        \node[labelnode] at (-0.2, 0.38) {\texttt{arr.xs}:};
        \foreach \i in {0,1,2,3} {
          \node[cell, fill=colorX, draw=colorXBorder, minimum width=1.3cm] (sx\i) at (\i*1.3 + 0.65, 0.38) {$x_\i$};
          \node[idx] at ($(sx\i.north) + (0, 0.18)$) {[\i]};
        }

        % ys
        \node[labelnode] at (-0.2, -0.38) {\texttt{arr.ys}:};
        \foreach \i in {0,1,2,3} {
          \node[cell, fill=colorY, draw=colorYBorder, minimum width=1.3cm] (sy\i) at (\i*1.3 + 0.65, -0.38) {$y_\i$};
          \node[idx] at ($(sy\i.south) + (0, -0.18)$) {[\i]};
        }
      \end{tikzpicture}
    \end{minipage}
    \caption{Structure-of-Arrays (SoA): Parallel contiguous buffers segregated by field.}
    \label{fig:layout-soa}
  \end{subfigure}

  \caption{Comparison of memory layouts alongside representative C++ container definitions for an array of 2D points $\langle x_i, y_i \rangle$. In (a) AoP, the array stores pointers to scattered heap tuples, requiring indirect memory loads. In (b) AoS, fields are stored inline contiguously in a single buffer. In (c) SoA, fields are decomposed into parallel contiguous buffers.}
  \label{fig:memory-layouts}
\end{figure}

\section{Implementation}
%% ---whole section outline --

%% visit all statements/types

%%   define a heuristic for the element type (e.g. tuple width), then find all
%%   references to arrays/vectors of that type

%%   recursively, deeply rewrite types of array/vector so long as there is an
%%   array/vector over a matched type. iterate until convergence to catch bugs
%%   that we'll describe. this step is for non-statement types (TODO: list them)

%%   rewrite all statements: prims are transformed in a container-/prim-specific
%%   way (which we outline below). statements and their attached types are
%%   transformed jointy. some are pass-through (type only), some need
%%   transformation

%%   unlike PreFlatten, there is no need ot iterate, we always converge after 1
%%   step

First, we will outline the high-level steps of the transformation (which are
identical for the two flattening strategies). Second, we will provide a detailed
description of the transformation that this pass applies to the types in the
program. Finally, we will provide the detailed \texttt{prim}-by-\texttt{prim}
transformation, \hl{specialized} for the AoS and SoA flavors of the
transformation

\subsection{High-level overview}

\hl{TODO: PROBABLY NEED A FORMAL GRAMMAR OF SSA TO MOTIVATE THIS}

\begin{enumerate}
\item Fix a heuristic that divides all SSA types into ``flattenable'' and
  ``un-flattenable.'' In this work, we define the ``tuple type width'' of any
  SSA type \texttt{t} as: if \texttt{t} is a tuple with $n$ members, its width
  its $n$, otherwise the width is 0.

\item Visit all non-\texttt{Statement} IR constructs (\texttt{datatype}
  constructors, function argument/return types, and block arguments) and
  iteratively flatten all flattenable types until convergence.

\item Visit all \texttt{Statement} IR constructs and iteratively flatten until
  convergence (i.e. flatten each \texttt{Statement}, perhaps emitting more
  \texttt{Statement}s, and then apply the flattening transformation to all
  emitted statements until they can no longer be further flattened).

\end{enumerate}

\subsection{Type transformation}
The core transformation is to rewrite an array-of-tuple to a tuple-of-array:
\begin{minted}{python}
  (a * b * c * ...) array ->
  a array * b array * c array * ...
\end{minted}
Since there exist constructs that convert the type \texttt{a array} to \texttt{a
  vector} (which is effectively a cast, for which MLTon inserts a copy in order
to prersve the property that \texttt{vector} objects are immutable after
construction), applying the optimization to \texttt{array} requires a either
matching transformation for \texttt{vector}, or a complicated re-un-flattening
in order to construct the contents of the un-flattened \texttt{vector} type.

In this document, we choose the former strategy, rewriting \texttt{vector} types
in addition to \texttt{array} types:
\begin{minted}{python}
  (a * b * c * ...) vector ->
  a vector * b vector * c vector * ...
\end{minted}

Likewise, it is necessary to traverse ``deeply'' inside of types, i.e. we need
to do the transformation:
\begin{minted}{python}
  ((a * b) array * c) ->
  ((a array * b array) * c)
\end{minted}
since, in order to construct an object of type \texttt{(a * b) array * c}, we
need an input object of type \texttt{(a * b) array}, which will be flattened by
the transformation below.

It is similarly necessary to apply the transformation iteratively, i.e. we need:
\begin{minted}{python}
  ((a * b) array * c) array ->
  ((a * b) array) array * c array ->
  (a array array * b array array) * c array
\end{minted}
since in order to construct the original type, we needed to have an object of
type \texttt{(a * b) array} in hand, and this object was transformed into an
\texttt{a array * b array} by the earlier action of the pass.

\subsection{Statement transformation}
The general form of an \texttt{SSA} statement is:
\begin{minted}{python}
  lhs: lhsTy = exp
\end{minted}
where \texttt{lhs} is a variable name, \texttt{lhsTy} is its type, and
\texttt{exp} is an expression.

The transformation on statements is divided into two cases: ``passthrough'' and
``rewrite.'' Every non-\texttt{PrimApp} statement is a ``passthrough''
statement. Taking \texttt{fTy} as the transformation on types above, these cases
are all
\begin{minted}{python}
  lhs: lhsTy = exp ->
  lhs: fTy(lhsTy) = exp
\end{minted}
where \texttt{exp} is some expression that contains no type expressions.

The general form of a \texttt{PrimApp} is:
\begin{minted}{python}
  lhs: lhsTy = prim PrimName[targ1, targ2, ...](arg1, arg2, ...) ->
\end{minted}
where \texttt{PrimName} is the name of a particular primitive operation (e.g.
\texttt{Word32\_add}, \texttt{Word32\_equal}, etc), \texttt{targ1, targ2, ...}
is a sequence of zero or more type arguments, and \texttt{arg1, arg2, ...} is a
sequence of zero or more variable names.


The ``passthrough'' case for \texttt{PrimApp} applies to all \texttt{PrimApp}s
that do not manipulate flattened containers. Its general form is:
\begin{minted}{python}
  lhs: lhsTy = prim PrimName[targ1, targ2, ...](arg1, arg2, ...) ->
  lhs: fTy(lhsTy) = prim PrimName[fTy(targ1), fTy(targ2), ...](arg1, arg2, ...)
\end{minted}
i.e.: leave all non-type portions of the expression unchanged, and transform all
type expressions according to the type transformation above.

The non-``passthrough,'' ``rewrite'' case applies to \texttt{PrimApp}s that
manipulate flattened containers. These transformations differ between the SoA
and AoS cases.

Beginning with the SoA case, the general form of the transformation is ``find
all containers of tuples, rewrite them to be tuples of containers.'' Beginning
with container-creation as an illustrative example, a single step of the
transformation is:
\begin{minted}{python}
  x: (a * b * ...) array = prim Array_alloc[a * b * ...](n) ->
  # Construct per-member arrays
  arr_a: a array = prim Array_alloc[a](n)
  arr_b: b array = prim Array_alloc[b](n)
  ...
  # Bind the tuple-of-arrays to the original name
  x: a array * b array * ... = tuple(arr_a, arr_b, ...)
\end{minted}
and now the other transformations of the pass are responsible for ensuring that
other uses of \texttt{x} in the program use the updated type (e.g. any function
signatures must be rewritten according to \texttt{tTy}, etc).

\hl{TODO: WE REALLY DO ONE MORE ITERATION AFTER THIS (i.e. first update
  statements until convergence, then update types, then run again). EXPLAIN?
  MIGHT BE WRONG}

Like as in the case of types, it is necessary to apply the transformation
iteratively until convergence in order to correctly handle the case of nested
types, e.g. after one round, we get:
\begin{minted}{python}
  arr: ((a * b) * c) array = prim Array_alloc[(a * b) * c](n) ->
  arr_ab: (a * b) array = prim Array_alloc[a * b]
  arr_c: c array = prim Array_alloc[a * b]
  x: ((a * b) array * c array) = tuple(arr_ab, arr_c)
\end{minted}
Now the type of \texttt{x} still admits flattening: this is a problem, since, as
we described above, the type transformation will eliminate all ocurrences of
flattenable types in the program. Hence any users of \texttt{a} will fail to
typecheck, and the resulting program is ill-formed. Applying another round of
``rewrite''-flattening to each of the resulting expressions resolves the issue:
\begin{minted}{python}
  arr_ab: (a * b) array = prim Array_alloc[a * b](n)
  arr_c: c array = prim Array_alloc[a * b](n)
  x: ((a * b) array * c array) = tuple(arr_ab, arr_c)
  ->
  arr_a: a array = prim Array_alloc[a](n)
  arr_b: b array = prim Array_alloc[b](n)
  arr_ab: a array * b array = tuple(arr_a, arr_b)
  arr_c: c array = prim Array_alloc[a * b]
  x: ((a array * b array) * c array) = tuple(arr_ab, arr_c)
\end{minted}
Most other \texttt{Array\_} transformations follow in the obvious way: \hl{TODO: CAN RETHINK ORDER}
\begin{enumerate}
\item \texttt{Array\_length} picks an arbitrary tuple member to pass the call on to)
\begin{minted}{python}
  arr: (a * b * ...) array = ...
  n: int = Array_length[a * b * ...](arr)
  -->
  # Earlier action of the pass means that now we have:
  #   arr: (a array * b array * ...) = ...
  arr_a = select (arr, 0)
  n: int = Array_length[a](arr_a)
\end{minted}

\item \texttt{Array\_sub} unpacks the tuple-of-arrays, reads out the necessary
  elements, and constructs the final output tuple:

\begin{minted}{python}
  arr: (a * b * ...) array = ...
  # Read element i
  x: (a * b * ...) = Array_sub[a * b * ...](arr, i)
  -->
  # Earlier action of the pass means that now we have:
  #   arr: (a array * b array * ...) = ...
  arr_a: a array = select (arr, 0)
  x_a: a = Array_sub[a](arr_a, i)
  arr_b: b array = select (arr, 1)
  x_b: b = Array_sub[b](arr_b, i)
  ...
  x: (a * b * ...) = tuple(x_a, x_b, ...)
\end{minted}

\item \texttt{Array\_update} unpacks the tuple-of-arrays and writes in the
  necessary elements:

\begin{minted}{python}
  arr: (a * b * ...) array = ...
  x: (a * b * ...) = ...
  _ = Array_update[a * b * ...](arr, i, x)
  -->
  # Earlier action of the pass means that now we have:
  #   arr: (a array * b array * ...) = ...
  arr_a: a array = select (arr, 0)
  x_a: a = select (x, 0)
  _ = Array_update[a](arr_a, i, x_a)
  arr_b: b array = select (arr, 1)
  x_b: b = select (x, 1)
  _ = Array_update[b](arr_b, i, x_b)
  ...
\end{minted}

\item \texttt{Array\_uninitIsNop} (which checks to see if raw data is dangerous
  for the garbage collector for the specified array type) is hardcoded to
  \texttt{false} for simplicity:

\begin{minted}{python}
  arr: (a * b * ...) array = ...
  isNop: bool = Array_uninitIsNop[(a * b * ...) array](arr)
  -->
  isNop: bool = false
\end{minted}

\item \texttt{Array\_uninit} (which fills an array with garbage data that is
  safe for the garbage collector) is passed on to the constituent members:

\begin{minted}{python}
  arr: (a * b * ...) array = ...
  _ = Array_uninit[a * b * ...](arr, n)
  -->
  # Earlier action of the pass means that now we have:
  #   arr: (a array * b array * ...) = ...
  arr_a: a array = select (arr, 0)
  _ = Array_uninit[a](arr_a, n)
  arr_b: b array = select (arr, 1)
  _ = Array_uninit[b](arr_b, n)
  ...
\end{minted}

\item \texttt{Array\_toArray} (which marks an array as initialized, and so safe
  for the garbage collector) is also passed on to the constituent members:

\begin{minted}{python}
  arr: (a * b * ...) array = ...
  arr': (a * b * ...) array = Array_toArray[a * b * ...](arr)
  -->
  # Earlier action of the pass means that now we have:
  #   arr: (a array * b array * ...) = ...
  arr_a: a array = select (arr, 0)
  arr'_a: a array = Array_toArray[a](arr_a)
  arr_b: b array = select (arr, 1)
  arr'_b: b array = Array_toArray[b](arr_b)
  ...
  arr': a array * b array * ... = tuple(arr'_a, arr'_b, ...)
\end{minted}

\end{enumerate}

The transformation additionally handles the \texttt{array -> vector}
transformation primitive (a copy into an immutable container); i.e.:

\begin{minted}{python}
  arr: (a * b * ...) array = ...
  vec: (a * b * ...) vector = Array_toVector[a * b * ...](arr)
  -->
  # Earlier action of the pass means that now we have:
  #   arr: (a array * b array * ...) = ...
  arr_a: a array = select (arr, 0)
  vec_a: a vector = Array_toVector[a](arr_a)
  arr_b: b array = select (arr, 1)
  vec_b: b vector = Array_toVector[b](arr_b)
  ...
  vec: a vector * b vector * ... = tuple(vec_a, vec_b, ...)
\end{minted}

Consequently, it is necessary to handle the full set of operations on
\texttt{vector}, i.e.:
\begin{enumerate}
\item \texttt{Vector\_length}: like \texttt{Array\_length}, arbitrarily forwards
  the call to the first tuple element:

\begin{minted}{python}
  vec: (a * b * ...) vector = ...
  n: int = Vector_length[a * b * ...](vec)
  -->
  # Earlier action of the pass means that now we have:
  #   vec: (a vector * b vector * ...) = ...
  vec_a = select (vec, 0)
  n: int = Vector_length[a](vec_a)
\end{minted}

\item \texttt{Vector\_sub}: analogous to \texttt{Array\_sub}

\begin{minted}{python}
  vec: (a * b * ...) vector = ...
  # Read element i
  x: (a * b * ...) = Vector_sub[a * b * ...](vec, i)
  -->
  # Earlier action of the pass means that now we have:
  #   vec: (a vector * b vector * ...) = ...
  vec_a: a vector = select (vec, 0)
  x_a: a = Vector_sub[a](vec_a, i)
  vec_b: b vector = select (vec, 1)
  x_b: b = Vector_sub[b](vec_b, i)
  ...
  x: (a * b * ...) = tuple(x_a, x_b, ...)
\end{minted}

\end{enumerate}

For the AoS version, we ``flatten'' by reinterpreting the AoP data as consective
elements in a larger array. Taking \texttt{Array\_alloc} again as a
representative example, this means:
\begin{minted}{python}
  x: (a * a * ...) array = prim Array_alloc[a * a * ...](n)
  -->
  # Scale up the array by the \# of tuple elements
  tupleSize: int = tupleWidth(a * a * ...)
  n': int = n * tupleWidth
  # Bind `x` to a flat array of the scaled-up size
  x: a array = Array_alloc[a](n')
\end{minted}

The other operations follow in the obvious way:
\begin{enumerate}
\item \texttt{Array\_length} scales down the length of the flattened array

\begin{minted}{python}
  arr: (a * a * ...) array = ...
  n: int = Array_length[a * a * ...](arr)
  -->
  # Earlier action of the pass means that now we have:
  #   arr: a array = ...
  tupleSize: int = tupleWidth(a * a * ...)
  n': int = Array_length[a](arr)
  n: int = n' / tupleSize
\end{minted}

\item \texttt{Array\_sub} reads out a slice from the flattened array and packs
  it into a tuple:

\begin{minted}{python}
  arr: (a * a * ...) array = ...
  # Read element i
  x: (a * a * ...) = Array_sub[a * a * ...](arr, i)
  -->
  # Earlier action of the pass means that now we have:
  #   arr: a array = ...
  tupleSize: int = tupleWidth(a * a * ...)
  x_0: a = Array_sub[a](arr, i * tupleSize)
  x_1: a = Array_sub[a](arr, i * tupleSize + 1)
  ...
  x_n: a = Array_sub[a](arr, i * tupleSize + tupleSize - 1)
  x: (a * a * ...) = tuple(x_0, x_1, ..., x_n)
\end{minted}

\item \texttt{Array\_sub} unpacks the input tuple and writes it into a slice:

\begin{minted}{python}
  arr: (a * a * ...) array = ...
  x: (a * a * ...) = ...
  _ = Array_update[a * a * ...](arr, i, x)
  -->
  # Earlier action of the pass means that now we have:
  #   arr: a array = ...
  tupleSize: int = tupleWidth(a * a * ...)
  x_0: a = select(x, 0)
  x_1: a = select(x, 1)
  ...
  _ = Array_update[a](arr, i * tupleSize + 0, x_0)
  _ = Array_update[a](arr, i * tupleSize + 1, x_1)
  ...
\end{minted}

\item \texttt{Array\_uninitIsNop} is hardcoded to \texttt{false}

\begin{minted}{python}
  arr: (a * a * ...) array = ...
  isNop: bool = Array_uninitIsNop[(a * a * ...) array](arr)
  -->
  isNop: bool = false
\end{minted}

\item \texttt{Array\_uninit} is passed on to the slice of the flattened array
  corresponding to the provided index:

\begin{minted}{python}
  arr: (a * a * ...) array = ...
  _ = Array_uninit[a * a * ...](arr, n)
  -->
  # Earlier action of the pass means that now we have:
  #   arr: a array = ...
  tupleSize: int = tupleWidth(a * a * ...)
  _ = Array_uninit[a](arr, n * tupleSize)
  _ = Array_uninit[a](arr, n * tupleSize + 1)
  ...
  _ = Array_uninit[a](arr, n * tupleSize + tupleSize - 1)
\end{minted}

\item \texttt{Array\_toVector} is passed on to the flattned array:

\begin{minted}{python}
  arr: (a * a * ...) array = ...
  vec: (a * a * ...) vector = Array_toVector[a * a * ...](arr)
  -->
  # Earlier action of the pass means that now we have:
  #   arr: a array = ...
  vec: a vector = Array_toVector[a](arr)
\end{minted}

\item \texttt{Vector\_sub} works analogously to \texttt{Array\_sub}

\begin{minted}{python}
  vec: (a * a * ...) vector = ...
  # Read element i
  x: (a * a * ...) = Vector_sub[a * a * ...](vec, i)
  -->
  # Earlier action of the pass means that now we have:
  #   vec: a vector = ...
  tupleSize: int = tupleWidth(a * a * ...)
  x_0: a = Vector_sub[a](vec, i * tupleSize)
  x_1: a = Vector_sub[a](vec, i * tupleSize + 1)
  ...
  x_n: a = Vector_sub[a](vec, i * tupleSize + tupleSize - 1)
  x: (a * a * ...) = tuple(x_0, x_1, ..., x_n)
\end{minted}

\end{enumerate}

\hl{TODO: DISCUSS SECOND CALL TO PROPAGATE TYPES?}

\section{Limitations}
As outlined above, this implementation only supports the case when all AoP
containers of a given type are flattened -- we cannot support a mix of flattened
and unflattened AoP containers with a given type using the strategy above. The
implementation discussed in this document simply emits an error if it encounters
an un-flattenable construct, none of which appear in the benchmark suites that
we evaluated this optimization against.

It is debatable whether supporting a fallback is required for the transformation
to be usable by real developers. The problem to worry about is that the
AoP->AoS/SoA transformation breaks pointer equality for array elements, i.e. for
an SML program like:
\begin{minted}{sml}
  val arr : (int * int) array = Array.array (3, (0, 0))
  val element = (1, 2)
  val _ = Array.update (arr, 1, element)
  val element' = Array.sub (arr, 1)
\end{minted}
Are \texttt{element} and \texttt{element'} physically identical? If this element
is \texttt{ShallowFlatten}ed, the answer is ``no:'' we cannot preserve pointer
equality, since the write/read result in pack/unpack operations that construct a
new tuple. From the perspective of strict StandardML as defined in the language
specification \citet{sml-def}, this is a non-issue: there is no way to observe
pointer equality of tuples in the source language, and so the compiler has
complete freedom. The only spec-compliant way to write this test would be to
wrap the element type in \texttt{ref}:
\begin{minted}{sml}
  val arr : (int * int) ref array = Array.array (3, ref (0, 0))
  val element = ref (1, 2)
  val _ = Array.update (arr, 1, element)
  val element' = Array.sub (arr, 1)
\end{minted}
Now we are guaranteed that \texttt{element = element'} holds, which is preserved
by \texttt{PreFlatten}.

However, both MLton and MPL extend the source language with constructs that
allow user programs to observe this behavior: both have \texttt{MLton.eq}
(physical equality for arbitrary values), and MPL additionally offers
compare-and-swap primitives to support implementing lock-free parallel
algorithms.

Since physical equality of tuples is not defined in \citet{sml-def}, one option
to resolve the issue would be to simply allow \texttt{MLton.eq}/compare-and-swap
to return \texttt{false} in these cases. Another simple option would to be
disable \texttt{ShallowFlatten} for all values of any types that participate in
one of the forbidden comparisons (i.e. a \texttt{MLton.eq} applied to an
\texttt{int * int} tuple anywhere in the program disables
\texttt{ShallowFlatten} for all containers of \texttt{int * int}), perhaps
refined by a dataflow analysis to determine if the elements in the forbidden
comparison may have originated in a container.

This results in a potentially-unpleasant performance cliff for users, but there
are various simple workarounds -- e.g. users can wrap any objects that must be
tested for pointer equality in \texttt{ref} cells (which are not impacted by
\texttt{ShallowFlatten}).

%% \hl{TODO}


\section{Comparison to more-agressive \texttt{DeepFlatten}}

\hl{TODO: implement an benchmark a version of deepflatten that kicks in for this
  case (i.e. no restriction on nested immutable types) / determine why it
  doesn't work}

